\section{DISCUSSION}
Using occupancy logs and field-derived PCMs under a shared setpoint-selection rule, this study shows that the value of temporal resolution depends on $K$, daily mean occupancy, and UTR.
These metrics can identify where real-time tracking may justify its cost.

\subsection{Why aggregation resolution changes the group-size result}
Prior simulations showed that the benefits of PCM integration diminish as more occupants share a zone, particularly beyond six occupants \parencite{jung_energy_2020}; at larger group scales, increasing group size and internal preference diversity likewise reduced the potential maximum satisfaction rate, while larger groups became less thermally sensitive \parencite{wang_enhancing_2026}.
The SGCM here follows the same directional pattern, crossing zero at $K\approx6$--8.
DGCM and RT-GCM, however, retained positive improvement at the same $K$ by excluding absent members.
This study therefore extends prior group-size findings by showing that comfort-aggregation outcomes depend not only on potential group size but also on realized occupancy (Figure~\ref{fig:OccuDensity}) and dynamic changes in occupant composition.
RT-GCM also exceeded DGCM by 6--8 percentage points after the initial rise with $K$, reflecting the value of knowing who is present now rather than who attended that day.
These predicted gains assume immediate setpoint selection and remain conditional on sensing reliability, HVAC response, and occupant adaptation.
% This is not proof that every real-time system will deliver the same field improvement: the metric is a predicted mean ``No change'' probability under a shared setpoint with immediate setpoint selection, so it identifies a theoretical value of temporal resolution whose actual delivery depends on sensing reliability, response, and adaptation.

\subsection{From comfort information to actionable control}
A high-resolution comfort representation is useful only when the HVAC system can respond to the resulting setpoint differences, reflecting the mismatch between comfort-model and control resolution reported by \textcite{ono_effects_2022}.
Although the simulation used a 0.1$^\circ$C grid, typical changed steps remained near 0.8$^\circ$C at higher mean occupancy and exceeded the 0.5$^\circ$C actionable threshold.
This indicates actuation-relevant potential rather than delivered field benefit; reviews report limited field implementation and persistent data, computation, BAS-integration, and scalability barriers \parencite{xie_review_2020,soleimanijavid_challenges_2024}.
UTR complements mean occupancy by indicating whether represented comfort profiles are being replaced, extending prior discussions of dynamic diversity in comfort-driven control \parencite{jung_energy_2020}.
The near-overlap across $K$ in Figure~\ref{fig:turnover_action}(a) shows that turnover is a better guide to update frequency than $K$ alone.
The three aggregation levels can consequently be treated as a deployment hierarchy: SGCM is defensible in stable assigned-seat zones, where UTR remained near 0.2--0.4 through office hours and the regular-member group matched presence; DGCM is a lower-complexity compromise when daily attendance varies but within-day composition stays relatively stable.
RT-GCM is most valuable where UTR instead fluctuated dynamically, with sparse attendance and frequent movement.
UTR can gate real-time updates to high-turnover periods, with deadbands or minimum hold times limiting excessive cycling.
Such gating does not remove sensing trade-offs in accuracy, latency, energy use, cost, scalability, and privacy \parencite{zafari_survey_2019}.
Because activity-based working outcomes depend on implementation and occupant use \parencite{marzban_review_2023}, these policy thresholds should be calibrated to each workplace.
% Even there, UTR can gate real-time updates to high-turnover periods, with deadbands or minimum hold times preventing excessive cycling.

\subsection{Limitations and research needs}
The results are conditional on one Singapore office building, four weeks, four rooms, and 39 PCMs; random PCM assignment broadened the combinations, but comfort remained predicted rather than closed-loop.
Energy use, HVAC dynamics and integration cost were not evaluated, so the deployment hierarchy should be interpreted as a comfort-based screening guide.
Mean comfort probability also cannot show whether lower aggregate comfort reflects minority disadvantage or a low attainable maximum under preference diversity; future work should evaluate occupant-level fairness and group-relative performance \parencite{wang_enhancing_2026}.
